\section{The metric tensor and arch length}

We already saw in the \textit{Tensor for beginners} notes, that the metric tensor is used to compute 
the length of a vector, the dot product between two vectors and the angle in between, and we also saw that 
it transforms using 2 covariant rules, hence is a (0,2)-tensor (bilinear form).

\begin{align*}
    g: V \times V \to \mathbb{R}\\
    g \left(\begingroup\color{orange}\vec{v}\endgroup, \begingroup\color{violet}\vec{w}\endgroup\right) \mapsto  \begingroup\color{orange}v^i\endgroup \begingroup\color{violet}w^j\endgroup \begingroup\color{blue}g_{ij}\endgroup
\end{align*}

\begin{align*}
    \begingroup\color{red}\tilde{g_{ij}}\endgroup &= F^k_i F^l_j \begingroup\color{blue}g_{kl}\endgroup\\
    \begingroup\color{blue}g_{kl}\endgroup &= B^i_k B^j_l \begingroup\color{red}\tilde{g_{ij}}\endgroup
\end{align*}

We'll here focus on using the metric tensor to compute lengths of curves in flat space.\\
We recall that in order to get the length of a curve in multi-variable calculus, basically we approximate 
the curve by a series of straight line segments, sum them up and take the limit as the $\Delta$ parameter used
to parametrize the curve approaches to 0.\\
At the end of the day, the key to get the curve length is getting the magnitude of the vector $\frac{dR}{d\lambda}$ if $\lambda$ is the parameter, and $R$ the 
position vector used to describe a point on the curve. Which is nothing but the tangent velocity vector along the curve:

\begin{align*}
    \text{arc length} = \int \norm{\frac{d\begingroup\color{violet}R\endgroup}{d\begingroup\color{brown}\lambda\endgroup}}d\begingroup\color{brown}\lambda\endgroup 
\end{align*}

Now, to get this magnitude, we can just compute the dot product of this vector with itself, as we know, and do that in cartesian coordinates for example:

\begin{align*}
    \norm{\frac{d\begingroup\color{violet}R\endgroup}{d\begingroup\color{brown}\lambda\endgroup}}^2 &=
    \frac{d\begingroup\color{violet}R\endgroup}{d\begingroup\color{brown}\lambda\endgroup} \cdot \frac{d\begingroup\color{violet}R\endgroup}{d\begingroup\color{brown}\lambda\endgroup}\\
    &= \left(\frac{\partial \begingroup\color{violet}R\endgroup}{\partial \begingroup\color{cyan}x\endgroup} \frac{d \begingroup\color{cyan}x\endgroup}{d \begingroup\color{brown}\lambda\endgroup} + 
    \frac{\partial \begingroup\color{violet}R\endgroup}{\partial \begingroup\color{cyan}y\endgroup} \frac{d \begingroup\color{cyan}y\endgroup}{d \begingroup\color{brown}\lambda\endgroup}\right) \cdot
    \left(\frac{\partial \begingroup\color{violet}R\endgroup}{\partial \begingroup\color{cyan}x\endgroup} \frac{d \begingroup\color{cyan}x\endgroup}{d \begingroup\color{brown}\lambda\endgroup} + 
    \frac{\partial \begingroup\color{violet}R\endgroup}{\partial \begingroup\color{cyan}y\endgroup} \frac{d \begingroup\color{cyan}y\endgroup}{d \begingroup\color{brown}\lambda\endgroup}\right)\\
    &= \left(\frac{d \begingroup\color{cyan}x\endgroup}{d \begingroup\color{brown}\lambda\endgroup}\right)^2 \left(\frac{\partial \begingroup\color{violet}R\endgroup}{\partial \begingroup\color{cyan}x\endgroup} \cdot \frac{\partial \begingroup\color{violet}R\endgroup}{\partial \begingroup\color{cyan}x\endgroup}\right) +
    2 \frac{d \begingroup\color{cyan}x\endgroup}{d \begingroup\color{brown}\lambda\endgroup} \frac{d \begingroup\color{cyan}y\endgroup}{d \begingroup\color{brown}\lambda\endgroup} \left(\frac{\partial \begingroup\color{violet}R\endgroup}{\partial \begingroup\color{cyan}x\endgroup} \cdot \frac{\partial \begingroup\color{violet}R\endgroup}{\partial \begingroup\color{cyan}y\endgroup}\right) +
    \left(\frac{d \begingroup\color{cyan}y\endgroup}{d \begingroup\color{brown}\lambda\endgroup}\right)^2  \left(\frac{\partial \begingroup\color{violet}R\endgroup}{\partial \begingroup\color{cyan}y\endgroup} \cdot \frac{\partial \begingroup\color{violet}R\endgroup}{\partial \begingroup\color{cyan}y\endgroup}\right)
\end{align*}

This ugly long product chain will look similar in other coordinate systems, like polar:

\begin{align*}
    &= \left(\frac{d \begingroup\color{magenta}r\endgroup}{d \begingroup\color{brown}\lambda\endgroup}\right)^2 \left(\frac{\partial \begingroup\color{violet}R\endgroup}{\partial \begingroup\color{magenta}r\endgroup} \cdot \frac{\partial \begingroup\color{violet}R\endgroup}{\partial \begingroup\color{magenta}r\endgroup}\right) +
    2 \frac{d \begingroup\color{magenta}r\endgroup}{d \begingroup\color{brown}\lambda\endgroup} \frac{d \begingroup\color{magenta}\theta\endgroup}{d \begingroup\color{brown}\lambda\endgroup} \left(\frac{\partial \begingroup\color{violet}R\endgroup}{\partial \begingroup\color{magenta}r\endgroup} \cdot \frac{\partial \begingroup\color{violet}R\endgroup}{\partial \begingroup\color{magenta}\theta\endgroup}\right) +
    \left(\frac{d \begingroup\color{magenta}\theta\endgroup}{d \begingroup\color{brown}\lambda\endgroup}\right)^2  \left(\frac{\partial \begingroup\color{violet}R\endgroup}{\partial \begingroup\color{magenta}\theta\endgroup} \cdot \frac{\partial \begingroup\color{violet}R\endgroup}{\partial \begingroup\color{magenta}\theta\endgroup}\right)
\end{align*}

If you carefully look at this, you will find out that it's very similar to when we computed the magnitude of a vector in our
\textit{Tensor for beginners} notes:

\begin{align*}
    \|\begingroup\color{orange}\vec{v}\endgroup\|^2 &= \begingroup\color{orange}\vec{v}\endgroup \cdot \begingroup\color{orange}\vec{v}\endgroup\\
    &= \left(\begingroup\color{cyan}v^1\endgroup \begingroup\color{blue}\vec{e_1}\endgroup + \begingroup\color{cyan}v^2\endgroup \begingroup\color{blue}\vec{e_2}\endgroup\right)
    \cdot \left(\begingroup\color{cyan}v^1\endgroup \begingroup\color{blue}\vec{e_1}\endgroup + \begingroup\color{cyan}v^2\endgroup \begingroup\color{blue}\vec{e_2}\endgroup\right)\\
    &= \bigl(\begingroup\color{cyan}v^1\endgroup\bigr)^2\left(\begingroup\color{blue}\vec{e_1}\endgroup\cdot\begingroup\color{blue}\vec{e_1}\endgroup\right)
    + 2 \begingroup\color{cyan}v^1v^2\endgroup\left(\begingroup\color{blue}\vec{e_1}\endgroup\cdot\begingroup\color{blue}\vec{e_2}\endgroup\right)
    + \bigl(\begingroup\color{cyan}v^2\endgroup\bigr)^2\left(\begingroup\color{blue}\vec{e_2}\endgroup\cdot\begingroup\color{blue}\vec{e_2}\endgroup\right) \\ \\
    &= \bigl(\begingroup\color{magenta}\tilde{v^1}\endgroup\bigr)^2\left(\begingroup\color{red}\tilde{\vec{e_1}}\endgroup\cdot\begingroup\color{red}\tilde{\vec{e_1}}\endgroup\right)
    + 2 \begingroup\color{magenta}\tilde{v^1v^2}\endgroup\left(\begingroup\color{red}\tilde{\vec{e_1}}\endgroup\cdot\begingroup\color{red}\tilde{\vec{e_2}}\endgroup\right)
    + \bigl(\begingroup\color{magenta}\tilde{v^2}\endgroup\bigr)^2\left(\begingroup\color{red}\tilde{\vec{e_2}}\endgroup\cdot\begingroup\color{red}\tilde{\vec{e_2}}\endgroup\right)
\end{align*}

The key is indeed getting those dot products like we did in those initial notes, and those are nothing more 
than the components of the metric tensor in different coordinate systems.\\

\begin{align*}
    g_{\color{cyan}c^i} = \begin{bmatrix}
        \frac{\partial \begingroup\color{violet}R\endgroup}{\partial \begingroup\color{cyan}x\endgroup} \cdot \frac{\partial \begingroup\color{violet}R\endgroup}{\partial \begingroup\color{cyan}x\endgroup} & \frac{\partial \begingroup\color{violet}R\endgroup}{\partial \begingroup\color{cyan}y\endgroup} \cdot \frac{\partial \begingroup\color{violet}R\endgroup}{\partial \begingroup\color{cyan}x\endgroup}\\
        \frac{\partial \begingroup\color{violet}R\endgroup}{\partial \begingroup\color{cyan}x\endgroup} \cdot \frac{\partial \begingroup\color{violet}R\endgroup}{\partial \begingroup\color{cyan}y\endgroup} & \frac{\partial \begingroup\color{violet}R\endgroup}{\partial \begingroup\color{cyan}y\endgroup} \cdot \frac{\partial \begingroup\color{violet}R\endgroup}{\partial \begingroup\color{cyan}y\endgroup}
    \end{bmatrix} &&
    g_{\color{magenta}p^i} = \begin{bmatrix}
        \frac{\partial \begingroup\color{violet}R\endgroup}{\partial \begingroup\color{magenta}r\endgroup} \cdot \frac{\partial \begingroup\color{violet}R\endgroup}{\partial \begingroup\color{magenta}r\endgroup} & \frac{\partial \begingroup\color{violet}R\endgroup}{\partial \begingroup\color{magenta}\theta\endgroup} \cdot \frac{\partial \begingroup\color{violet}R\endgroup}{\partial \begingroup\color{magenta}r\endgroup}\\
        \frac{\partial \begingroup\color{violet}R\endgroup}{\partial \begingroup\color{magenta}r\endgroup} \cdot \frac{\partial \begingroup\color{violet}R\endgroup}{\partial \begingroup\color{magenta}\theta\endgroup} & \frac{\partial \begingroup\color{violet}R\endgroup}{\partial \begingroup\color{magenta}\theta\endgroup} \cdot \frac{\partial \begingroup\color{violet}R\endgroup}{\partial \begingroup\color{magenta}\theta\endgroup}
    \end{bmatrix}
\end{align*}

Allright, now using the Einstein notation to save some writing:

\begin{align*}
    \norm{\frac{d\begingroup\color{violet}R\endgroup}{d\begingroup\color{brown}\lambda\endgroup}}^2 &=
    \frac{d\begingroup\color{violet}R\endgroup}{d\begingroup\color{brown}\lambda\endgroup} \cdot \frac{d\begingroup\color{violet}R\endgroup}{d\begingroup\color{brown}\lambda\endgroup}\\
    &= \frac{d \begingroup\color{cyan}c^i\endgroup}{d \begingroup\color{brown}\lambda\endgroup}\frac{d \begingroup\color{cyan}c^j\endgroup}{d \begingroup\color{brown}\lambda\endgroup}
    \left(\frac{\partial \begingroup\color{violet}R\endgroup}{\partial \begingroup\color{cyan}c^i\endgroup} \cdot \frac{\partial \begingroup\color{violet}R\endgroup}{\partial \begingroup\color{cyan}c^j\endgroup}\right) = \frac{d \begingroup\color{cyan}c^i\endgroup}{d \begingroup\color{brown}\lambda\endgroup}\frac{d \begingroup\color{cyan}c^j\endgroup}{d \begingroup\color{brown}\lambda\endgroup} \begingroup\color{blue}g_{ij}\endgroup\\
    &= \frac{d \begingroup\color{magenta}p^i\endgroup}{d \begingroup\color{brown}\lambda\endgroup}\frac{d \begingroup\color{magenta}p^j\endgroup}{d \begingroup\color{brown}\lambda\endgroup}
    \left(\frac{\partial \begingroup\color{violet}R\endgroup}{\partial \begingroup\color{magenta}p^i\endgroup} \cdot \frac{\partial \begingroup\color{violet}R\endgroup}{\partial \begingroup\color{magenta}p^j\endgroup}\right) = \frac{d \begingroup\color{magenta}p^i\endgroup}{d \begingroup\color{brown}\lambda\endgroup}\frac{d \begingroup\color{magenta}p^j\endgroup}{d \begingroup\color{brown}\lambda\endgroup} \begingroup\color{red}\tilde{g_{ij}}\endgroup
\end{align*}

And its transformation rules:

\begin{align*}
    \begingroup\color{red}\tilde{g_{ij}}\endgroup = \frac{\partial \begingroup\color{cyan}c^k\endgroup}{\partial \begingroup\color{magenta}p^i\endgroup} \frac{\partial \begingroup\color{cyan}c^l\endgroup}{\partial \begingroup\color{magenta}p^j\endgroup} \begingroup\color{blue}g_{kl}\endgroup \\
    \begingroup\color{blue}g_{kl}\endgroup = \frac{\partial \begingroup\color{magenta}p^i\endgroup}{\partial \begingroup\color{cyan}c^k\endgroup} \frac{\partial \begingroup\color{magenta}p^j\endgroup}{\partial \begingroup\color{cyan}c^l\endgroup} \begingroup\color{red}\tilde{g_{ij}}\endgroup
\end{align*}